\subsection{Pregunta 4: tolerancia a fallos y reconvergencia}
\autor{Benjamín}

\textbf{Escenario:} comunicación entre un host de la LAN de Raze (Paint Shells, \texttt{172.16.64.2}) y uno de la LAN de Reyna (Dismiss, \texttt{183.112.0.2}). Se suspende la interfaz del enlace 5.0.0.0 (Raze--Neon) y se comprueba el camino alternativo por el enlace 6.0.0.0 (Killjoy--Fade). Al terminar se restaura el enlace. \pend{confirmar las direcciones de los hosts con \cmd{ipconfig} si el criterio de subredes cambia}

\subsubsection*{Ruta original}

Desde Paint Shells se ejecutó \cmd{tracert 183.112.0.2}, y en Raze se consultó la ruta hacia la LAN de Reyna.

\evid{Q4_caida_enlace/ruta_original.png}{Traza de Paint Shells a Dismiss con todos los enlaces activos.}
\evid{Q4_caida_enlace/show_ip_route_raze_antes.png}{Tabla de rutas de Raze hacia la LAN de Reyna antes de la caída.}

\pend{ruta activa router por router, y el protocolo (B u O) con que Raze aprendió la ruta}

\subsubsection*{Caída del enlace 5.0.0.0}

Se lanzó un ping continuo desde Paint Shells (\cmd{ping -t 183.112.0.2}) y se suspendió la interfaz del enlace en Raze:

\begin{lstlisting}
Raze# configure terminal
Raze(config)# interface GigabitEthernet0/1
Raze(config-if)# shutdown
\end{lstlisting}

\evid{Q4_caida_enlace/ping_continuo_caida.png}{Ping continuo durante la caída: paquetes perdidos hasta que se restablece la comunicación.}

\subsubsection*{Nueva ruta y tiempo de convergencia}

\evid{Q4_caida_enlace/ruta_alternativa.png}{Traza tras suspender el enlace 5.0.0.0.}
\evid{Q4_caida_enlace/show_ip_route_raze_despues.png}{Tabla de rutas de Raze tras la caída.}

\pend{nueva ruta, número de paquetes perdidos y tiempo de convergencia observado (cada ping sin respuesta equivale a cerca de un segundo), y justificación: detección inmediata de la caída de interfaz, recálculo en OSPF dentro del AS 10 y selección de otra ruta en BGP/OSPF}

\subsubsection*{Restauración}

\begin{lstlisting}
Raze(config-if)# no shutdown
\end{lstlisting}

\pend{confirmar con ping y traza que el tráfico volvió a la ruta original, y adjuntar la captura}
